\section{ACSD Extension (Prospective)}
\label{sec:acsd}

The current empirical study identifies a practical optimum at moderate block
size ($k{=}8$). Beyond this point, accepted block length can increase while net
speedup declines, indicating that fixed-policy drafting is vulnerable to
acceptance collapse and long-tail latency. This motivates a targeted extension:
Adaptive Cascaded Speculative Decoding (ACSD), which keeps the standard
verifier but adapts policy online.

\subsection{Design objective}

Let speculative-cycle latency be
\begin{equation}
  L = \frac{T_{\text{draft}} + T_{\text{verify}}}{\tau},
\end{equation}
where $\tau$ is expected accepted tokens per cycle. In fixed-policy speculative
decoding, poorly matched samples can reduce $\tau$ and increase total verify
steps, producing heavy-tail runtime. ACSD explicitly targets this effect by
adapting draft policy at sample level.

\subsection{Proposed approach}

ACSD combines three online controls while preserving standard target-side
verification:
\begin{enumerate}
  \item \textbf{Adaptive $k$:} choose $k\in\{4,8,16\}$ from rolling acceptance,
        using larger blocks only when acceptance supports it.
  \item \textbf{Cascaded draft switch:} use 0.5B as a fast primary drafter and
        switch to 1.5B rescue drafting when acceptance remains low for
        consecutive verify steps.
  \item \textbf{Tail guardrail:} if acceptance stays below a threshold after a
        minimum generated length, fall back to target-only autoregressive
        decoding for the remainder of that sample.
\end{enumerate}

This design preserves verifier compatibility and allows direct comparison with
vanilla speculative decoding under the same manifests, prompts, metrics, and
hardware.

\begin{figure}[t]
\centering
\resizebox{0.98\columnwidth}{!}{%
\begin{tikzpicture}[
      node distance=4.5mm and 5mm,
      box/.style={rectangle, draw, rounded corners=1mm, align=center, minimum height=6mm, text width=30mm, inner sep=1.5mm},
      decision/.style={diamond, draw, align=center, aspect=2, inner sep=1.1mm, text width=22mm},
      line/.style={-Latex, thick, shorten >=1pt, shorten <=1pt}
]
\node[box] (start) {Start sample\\init primary draft, $k{=}8$};
\node[box, below=of start] (draft) {Draft proposes $k$ tokens\\(0.5B primary or 1.5B rescue)};
\node[box, below=of draft] (verify) {Target verifies draft block\\compute rolling $\alpha$};
\node[decision, below=of verify] (high) {$\alpha \ge \tau_{\text{high}}$?};
\node[box, left=of high, xshift=-8mm] (kup) {Increase $k$\\toward 16};
\node[decision, below=of high] (low) {$\alpha \le \tau_{\text{low}}$?};
\node[box, left=of low, xshift=-8mm] (kdown) {Reduce $k$\\toward 4};
\node[decision, below=of low] (rescue) {Low-accept streak\\$\ge c$?};
\node[box, left=of rescue, xshift=-8mm] (switch) {Switch to 1.5B rescue\\for $h$ cycles};
\node[decision, below=of rescue] (fallback) {$\alpha < \tau_{\text{AR}}$ and\\length $\ge L_{\min}$?};
\node[box, right=of fallback, xshift=8mm] (ar) {Fallback to target-only\\AR for tail};
\node[box, below=of fallback] (loop) {Continue next ACSD cycle};
\node[box, below=of loop] (end) {EOS or max tokens\\end sample};
\coordinate (loopRightTop) at ([xshift=18mm]draft.east);
\coordinate (loopRightBottom) at ([xshift=18mm]end.east);

\draw[line] (start) -- (draft);
\draw[line] (draft) -- (verify);
\draw[line] (verify) -- (high);
\draw[line] (high) -- node[above, sloped, font=\scriptsize] {yes} (kup);
\draw[line] (kup.east) to[out=0, in=180] (draft.west);
\draw[line] (high) -- node[right, font=\scriptsize] {no} (low);
\draw[line] (low) -- node[above, sloped, font=\scriptsize] {yes} (kdown);
\draw[line] (kdown.east) to[out=0, in=180] (draft.west);
\draw[line] (low) -- node[right, font=\scriptsize] {no} (rescue);
\draw[line] (rescue) -- node[above, sloped, font=\scriptsize] {yes} (switch);
\draw[line] (switch.east) to[out=0, in=180] (draft.west);
\draw[line] (rescue) -- node[right, font=\scriptsize] {no} (fallback);
\draw[line] (fallback) -- node[above, font=\scriptsize] {yes} (ar);
\draw[line] (fallback) -- node[right, font=\scriptsize] {no} (loop);
\draw[line] (ar) |- (end.east);
\draw[line] (loop) -- (end);
\draw[line] (end.east) -- (loopRightBottom) -- (loopRightTop) -- (draft.east);
\end{tikzpicture}%
}
\caption{Adaptive Cascaded Speculative Decoding (ACSD) control loop. ACSD
adapts block size from rolling acceptance, switches from the 0.5B primary to a
1.5B rescue drafter when acceptance remains low, and applies target-only
autoregressive fallback for pathological tails.}
\label{fig:acsd-flow}
\end{figure}

\subsection{Evaluation plan}

The extension will be evaluated as a strict incremental study rather than a new
protocol. We will retain:
\begin{itemize}
  \item the same target model and task manifests,
  \item the same deterministic and stochastic regimes,
  \item the same reporting metrics ($S$, $\alpha$, $B_{\text{eff}}$, quality deltas),
  \item the same artifact schema for per-configuration CSV outputs, with ACSD
        metadata fields for adaptive-$k$ usage, rescue usage, and fallback rate.
\end{itemize}

Primary comparison points are the current best fixed-policy autoregressive
settings ($k{=}8$ deterministic) and the high-$k$ settings ($k{=}16$) where
fixed-policy decoding showed diminishing returns.

\subsection{Status}

At the time of writing, ACSD is a planned extension and is not included in the
quantitative results of Section~\ref{sec:results}. We include it here to make
the research trajectory explicit and to separate completed evidence from
forward-looking engineering work.
